\section{Adaptive Terminology Memory}
\label{sec:adaptive-memory}

AutoTerm-SST maintains two different notions of terminology memory. The
\emph{open memory} is a large offline resource derived from Wikidata/Wikipedia
and curated glossaries. It is stored as manifest-driven snapshots with
precomputed MaxSim indexes. The \emph{working glossary} is the compact active
slice used during a live session. This distinction is important: a large memory
improves coverage, but injecting many broad references directly into every
prompt can reduce precision and make the UI harder to inspect.

\paragraph{Open memory snapshots.}
Each snapshot maps a preset id to term files, metadata, and a MaxSim index. The
runtime never stores large indexes in the code repository; instead, manifests
point to a runtime root on the demo host. Current presets include a common
10k-term slice, compact domain slices such as NLP and medicine, a curated
238-term ACL glossary, and broad open memories with 100k and 1M terms.

\paragraph{Audio-native active glossary router.}
The default router is \texttt{embedding\_refs}. It is not a trained topic
classifier and does not use ASR, source transcripts, generated target text, LLM
classification, or manual glossary terms for routing. Instead, it combines two
signals already available from speech-side retrieval:
\begin{equation}
\begin{split}
score(d) = & \lambda_e \cos(\mathrm{EMA}(q_{\mathrm{speech}}), c_d) \\
           & + \lambda_r \mathrm{vote}_d(\mathrm{references}) - \mathrm{ambiguity},
\end{split}
\end{equation}
where \(q_{\mathrm{speech}}\) is the pooled speech query embedding, \(c_d\) is a
domain centroid built offline from a slice's MaxSim text embeddings, and
reference votes come from metadata such as source preset and domain. The default
weights are \(\lambda_e=0.65\) and \(\lambda_r=0.35\).

\paragraph{Conservative switching.}
The router is a policy for selecting an active glossary, not an open-world topic
detector. It switches only after warmup, minimum update interval, confidence
threshold, top-vs-runner-up margin, repeated consistent windows, and target
index readiness. If evidence is weak or ambiguous, the session stays on the
common glossary or the current active slice. Cold index loads are non-blocking:
the agent preloads a target index in the background and activates it only after
the retrieval plugin reports that the index is ready.

\paragraph{Prompt and UI budgets.}
The retriever can return up to the UI evidence budget, and prompt injection uses
the same default budget:
\begin{center}
\small
\texttt{RASST\_PROMPT\_TOP\_K}=10\\
\texttt{RASST\_UI\_TOP\_K}=10
\end{center}
This cap is part of the system design. Broad memories may contain useful terms,
but the live prompt should remain a compact, inspectable working context.
